\documentclass[UTF8,a4paper,11pt,fontset=fandol]{ctexart}
\newcommand{\DocTitle}{第三次实验：逻辑回归与贷款违约预测}
\newcommand{\DocSubject}{2026年秋季课程B；135分钟；逻辑回归实验}
\newcommand{\DocShortTitle}{实验03：逻辑回归}
\newcommand{\DocType}{学生实验手册}
% Shared style for integrated experiment handouts.
\usepackage[left=2.15cm,right=2.15cm,top=2.25cm,bottom=2.15cm,headsep=0.7cm]{geometry}
\usepackage{amsmath,amssymb}
\usepackage{booktabs,array,tabularx,longtable,multirow}
\usepackage{enumitem}
\usepackage{xcolor}
\usepackage{graphicx}
\usepackage{tikz}
\usetikzlibrary{arrows.meta,positioning,calc,shapes.geometric}
\usepackage{tcolorbox}
\usepackage{listings}
\usepackage{hyperref}
\usepackage{fancyhdr}
\usepackage{chngcntr}

\definecolor{Ink}{HTML}{17233C}
\definecolor{Navy}{HTML}{173B68}
\definecolor{Blue}{HTML}{2878B5}
\definecolor{Cyan}{HTML}{2A9D8F}
\definecolor{Amber}{HTML}{E9A23B}
\definecolor{Red}{HTML}{C84B4B}
\definecolor{Paper}{HTML}{F6F8FB}
\definecolor{PaleBlue}{HTML}{EAF2FA}
\definecolor{PaleCyan}{HTML}{EAF7F4}
\definecolor{PaleAmber}{HTML}{FFF5E4}
\definecolor{Rule}{HTML}{D6DEE8}
\definecolor{Muted}{HTML}{5B677A}

\hypersetup{
  colorlinks=true,
  linkcolor=Navy,
  urlcolor=Blue,
  citecolor=Cyan,
  pdfauthor={《人工智能数学原理与算法》课程组},
  pdftitle={\DocTitle},
  pdfsubject={\DocSubject}
}

\ctexset{
  section={format=\Large\bfseries\color{Navy},beforeskip=1.4ex plus .2ex,afterskip=.8ex},
  subsection={format=\large\bfseries\color{Ink},beforeskip=1.2ex plus .2ex,afterskip=.55ex},
  subsubsection={format=\normalsize\bfseries\color{Blue},beforeskip=1ex,afterskip=.35ex}
}

\setlength{\parindent}{2em}
\setlength{\parskip}{0.22em}
\setlength{\emergencystretch}{2em}
\linespread{1.14}
\setlist[itemize]{leftmargin=2em,itemsep=.25em,topsep=.35em}
\setlist[enumerate]{leftmargin=2.2em,itemsep=.28em,topsep=.35em}
\renewcommand{\arraystretch}{1.30}
\newcolumntype{Y}{>{\raggedright\arraybackslash}X}
\newcolumntype{C}[1]{>{\centering\arraybackslash}p{#1}}
\newcolumntype{L}[1]{>{\raggedright\arraybackslash}p{#1}}

\newcommand{\R}{\mathbb{R}}
\newcommand{\E}{\mathbb{E}}
\newcommand{\MSE}{\operatorname{MSE}}
\newcommand{\RMSE}{\operatorname{RMSE}}
\newcommand{\argmin}{\operatorname*{arg\,min}}
\newcommand{\vect}[1]{\boldsymbol{#1}}
\newcommand{\mat}[1]{\mathbf{#1}}
\newcommand{\term}[1]{\textcolor{Blue}{\textbf{#1}}}
\newcommand{\code}[1]{\texttt{#1}}
\numberwithin{equation}{subsection}

\tcbset{boxrule=.7pt,arc=1.6mm,left=2.5mm,right=2.5mm,top=1.8mm,bottom=1.8mm,before skip=7pt,after skip=7pt}
\newtcolorbox{keybox}[1][核心结论]{colback=PaleBlue,colframe=Blue,title={#1},fonttitle=\bfseries,coltitle=white}
\newtcolorbox{examplebox}[1][贯穿例题]{colback=PaleCyan,colframe=Cyan,title={#1},fonttitle=\bfseries,coltitle=white}
\newtcolorbox{questionbox}[1][检查点]{colback=PaleAmber,colframe=Amber,colbacktitle=PaleAmber,title={#1},fonttitle=\bfseries,coltitle=Ink}
\newtcolorbox{pitfallbox}[1][易错提醒]{colback=red!3,colframe=Red,title={#1},fonttitle=\bfseries,coltitle=white}
\newtcolorbox{teacherbox}[1][教师提示]{colback=Paper,colframe=Rule,title={#1},fonttitle=\bfseries,coltitle=Muted}
\newtcolorbox{proofbox}[1][推导与证明]{colback=Blue!3,colframe=Navy,title={#1},fonttitle=\bfseries,coltitle=white}

\lstdefinestyle{python}{
  language=Python,
  basicstyle=\ttfamily\small,
  keywordstyle=\color{Blue}\bfseries,
  commentstyle=\color{Cyan!70!black},
  stringstyle=\color{Red!80!black},
  showstringspaces=false,
  columns=fullflexible,
  keepspaces=true,
  frame=single,
  rulecolor=\color{Rule},
  backgroundcolor=\color{Paper},
  breaklines=true,
  numbers=left,
  numberstyle=\tiny\color{Muted},
  xleftmargin=1.5em,
  framexleftmargin=1.2em,
  aboveskip=.7em,
  belowskip=.7em
}
\lstset{style=python}

\pagestyle{fancy}
\fancyhf{}
\fancyhead[L]{\small\color{Muted}人工智能数学原理与算法}
\fancyhead[R]{\small\color{Muted}\DocShortTitle}
\fancyfoot[L]{\small\color{Muted}\DocType}
\fancyfoot[R]{\small\color{Muted}\thepage}
\renewcommand{\headrulewidth}{0pt}
\renewcommand{\footrulewidth}{0pt}

\newcommand{\makeexperimenttitle}[4]{
  \hypersetup{pageanchor=false}
  \begin{titlepage}\thispagestyle{empty}
  \begin{tikzpicture}[remember picture,overlay]
    \fill[Ink] (current page.north west) rectangle ([yshift=-9.4cm]current page.north east);
    \fill[Blue] ([yshift=-9.4cm]current page.north west) rectangle ([yshift=-9.68cm]current page.north east);
    \fill[Cyan] ([xshift=2.15cm,yshift=-1.55cm]current page.north west) rectangle ([xshift=2.27cm,yshift=-8.55cm]current page.north west);
    \fill[Paper] ([xshift=2.15cm,yshift=2.0cm]current page.south west) rectangle ([xshift=-2.15cm,yshift=6.15cm]current page.south east);
  \end{tikzpicture}
  \vspace*{1.45cm}
  {\color{white}\Large 人工智能数学原理与算法\par}\vspace{.55cm}
  {\color{white}\fontsize{31}{39}\selectfont\bfseries #1\par}\vspace{.35cm}
  {\color{white}\fontsize{21}{28}\selectfont #2\par}
  \vfill
  \begin{tcolorbox}[colback=white,colframe=white,boxrule=0pt,arc=0mm,left=7mm,right=7mm,top=6mm,bottom=6mm,width=\textwidth]
    \begin{tabularx}{\linewidth}{L{2.6cm}Y}
      \textcolor{Muted}{材料类型} & \DocType\\
      \textcolor{Muted}{实验安排} & #3\\
      \textcolor{Muted}{贯穿案例} & #4\\
      \textcolor{Muted}{适用对象} & 已完成机器学习导论、具备高中数学与基础计算机操作能力的本科生\\
    \end{tabularx}
  \end{tcolorbox}
  \vspace{.7cm}{\color{Muted}\small 版本：2026年秋季学期整合稿\hfill 源文件与 PDF 同目录交付\par}
  \end{titlepage}
  \hypersetup{pageanchor=true}
  \pagenumbering{roman}\tableofcontents\clearpage\pagenumbering{arabic}
}

\usepackage{xurl,bookmark}
\setmonofont[BoldFont=DejaVuSansMono-Bold.ttf,ItalicFont=DejaVuSansMono-Oblique.ttf]{DejaVuSansMono.ttf}
\setCJKmonofont[AutoFakeBold=2]{FandolFang-Regular.otf}
\setlength{\headheight}{15pt}
\setcounter{tocdepth}{1}
\hypersetup{pdfauthor={钟志诚（主讲）；曾有成（整理）},pdftitle={第三次实验：分阶段实验手册}}
\lstset{numbers=none,xleftmargin=0pt,framexleftmargin=0pt,basicstyle=\ttfamily\small}
\newcommand{\file}[1]{\nolinkurl{#1}}
\begin{document}
\hypersetup{pageanchor=false}
\begin{titlepage}
\thispagestyle{empty}
\vspace*{1.4cm}
{\Large\color{Navy}人工智能数学原理与算法 B\par}
\vspace{1.0cm}
{\huge\bfseries\color{Ink}第三次实验\par}
\vspace{.4cm}
{\LARGE\bfseries\color{Navy}逻辑回归与贷款违约预测\par}
\vspace{.8cm}
{\Large 分阶段实验手册\par}
\vspace{1.1cm}
\begin{tabular}{@{}ll}
主讲 & 钟志诚\\[.5em]
整理 & 曾有成\\[.5em]
单位 & 中国科学技术大学\\
 & 人工智能与数据科学学院
\end{tabular}
\vspace{1cm}
\begin{keybox}[从数据到可解释的分类流程]
数据清洗与划分 $\rightarrow$ 数值表示 $\rightarrow$ 线性得分与概率
$\rightarrow$ 交叉熵与正则化 $\rightarrow$ 梯度更新 $\rightarrow$ 独立评估。
\end{keybox}
\vfill
2026 年秋季\quad 建议 3 学时（135 分钟）\\
C0--C7 分阶段完成；日期与提交平台以课程通知为准。\\
随包数据为天池真实贷款记录的固定课堂子集。
\end{titlepage}
\hypersetup{pageanchor=true}
\pagenumbering{roman}\tableofcontents\clearpage\pagenumbering{arabic}
\section{如何使用这份实验材料}
\subsection{课前准备与文件入口}
先阅读《实验03\_前置讲义》的第 1--6 节，课前约 30 分钟；课堂通过手算与代码继续学习第 7--10 节。实验总时长建议 135 分钟，不把课前安装时间算入课堂。

沿用实验 01 的课程环境。终端进入本实验的 \file{code/} 目录，运行：
\begin{lstlisting}[language=bash]
conda activate ustc-ai-experiments-26fall-B
python -c "import torch, pandas, matplotlib; print(torch.__version__)"
python checkpoints.py C0
\end{lstlisting}
若缺少模块，在已激活环境运行 \texttt{python -m pip install -r requirements.txt}。本实验在 CPU（Central Processing Unit，中央处理器）上运行即可，无需配置加速硬件。依赖文件列出本次验证过的版本，已有可用环境不必重复安装。
\par\noindent\begin{tabularx}{\linewidth}{L{4.2cm}Y}\toprule
文件 & 用途\\\midrule
\file{student.py} & 唯一主要补全文件，所有 TODO 均有 Hint\\
\file{checkpoints.py} & C0--C7 小样例检查；只检查指定阶段\\
\file{support.py}\newline\file{workflow.py} & 已提供的读表、日期编码、绘图、模型选择与测试流程\\
\file{run_lab.py} & 补全后运行整个实验的普通 Python 入口\\
\file{logic.ipynb} & 按 checkpoint 组织的 Notebook 入口，复用同一份学生代码\\
\file{data/classroom/} & 真实贷款课堂子集与划分记录，可离线运行\\
\file{report-template.md} & 实验记录模板，含每阶段证据与总结\\\bottomrule
\end{tabularx}
\subsection{每一阶段都按同一节奏进行}
阅读目标，先预测小例子的输出，再运行并改动一个值；补全指定函数；运行对应检查；把实现思路和一条数值验证留作报告素材，按第 7 节整理。\texttt{Ck PASS} 只表示所列小样例通过，不能代替解释与完整实验。

Notebook 通过普通 \texttt{import student} 调用 \file{student.py}。修改该文件后，重启内核并从头执行到当前阶段，避免用到旧函数；不要在两处维护不同答案。初始骨架中的 \texttt{return None} 是明确的待完成占位，不是算法返回值。
\begin{pitfallbox}[不要跳过形状与数据来源]
随包数据不是教材原始贷款数据。先完成 C0；C1--C7 未完成时出现断言失败或类型报错是预期行为，按 checkpoint 提示定位。不要把全部报错都归因于环境。
\end{pitfallbox}
\clearpage
\section{135 分钟课堂路线}
\subsection{教师活动、学生活动与可观察产出}
\par\noindent\begin{tabularx}{\linewidth}{L{1.8cm}L{2.2cm}Y Y}
\toprule 阶段/分钟 & 内容 & 教师活动 & 学生活动与产出\\\midrule
C0 / 10 & 环境、形状 & 连接前两次实验，演示矩阵乘法 & 运行检查，记录得分与形状\\
C1 / 15 & 读表、编码 & 说明标签和编号职责 & 补两个函数，列出缺失数和编码结果\\
C2 / 20 & 填充、尺度 & 手算均值与训练统计量 & 补三个函数，输出张量维度\\
C3 / 15 & 批次 & 用五条数据演示成对取样 & 补数据集方法和 loader，观察末批\\
C4 / 20 & 得分与概率 & 画出 sigmoid，解释参数对象 & 补模型，手算与代码核对\\
C5 / 20 & 损失、正则 & 推导交叉熵，讲解接口 & 补损失和惩罚，核对数值\\
C6 / 20 & 一次更新 & 板书梯度与五步训练循环 & 补训练，验证两次连续更新\\
C7 / 15 & 选择、评价 & 讲解四类计数和三种数据集 & 补指标，运行方案比较与退出卡\\\bottomrule
\end{tabularx}
\medskip
合计 $10+15+20+15+20+20+20+15=135$ 分钟。每段讲解后立即上机，C6/C7 完整记录可课后完善；课堂不临时增加必须完成的大范围调参。
\subsection{完成标准与依赖关系}
C0 独立可运行。C1、C2 的小测试彼此独立；C3 使用提供的小张量；C4 建立模型；C5 需要 C4；C6 需要 C3--C5；C7 指标小测试可独立完成，但完整实验需要全部阶段。

\textbf{必做：}C0--C7、课堂手算、默认三方案比较、冻结后最终测试、错误分析与实验反馈。\textbf{选做：}在看测试结果之前，只在训练/验证上改变学习率或选择阈值。若测试已揭晓，再修改方案仅算探索，不能把同一测试集当成新的独立证据。
\begin{teacherbox}[建议板书]
左侧：原始表、训练统计量和张量形状；中间：$\mathbf X\boldsymbol\theta\to\mathbf p\to R\to\nabla R$；右侧：训练/验证/测试职责与混淆矩阵。各记号定义见前置讲义第 3、5--9 节。
\end{teacherbox}
\begin{questionbox}[检查规则]
每个 checkpoint 在课堂自查时核对一条数值输出，并说清自己的解释。报告逐项记录关键实现、小例子的输入及预期与实际输出、结果解释；综合训练图表集中展示。未通过时记录具体函数、现象与当前理解。
\end{questionbox}
\clearpage
\section{C0--C2：从表格到张量}
\subsection{C0：先运行一个已完成的小例子}
运行 \texttt{python checkpoints.py C0}，预期显示课堂训练表 \texttt{(8000,47)}、缺失标签 0、矩阵乘法结果 \texttt{[[2.0],[-1.0]]}。原表保留 47 列，本次只选八个输入字段，编号与标签不进入模型。把例子中的权重 1 改为 2，先预测输出再核对；改回后继续。
\subsection{C1：缺失标签与就业年限}
在 \file{student.py} 补 \texttt{filter\_labels} 和 \texttt{encode\_employment}。前者只删标签缺失行，不在这一步填充输入；后者处理缺失值、单数 year、复数 years 和两个特殊档位。
\begin{lstlisting}[language=Python]
import pandas as pd
import support
train = support.load_split(support.DATA_DIR, "train")
print(train.head())
print(train.isna().sum())
print(train["isDefault"].value_counts(dropna=False))
\end{lstlisting}
运行 \texttt{python checkpoints.py C1}。应保留示例中特征缺失但标签有效的行；年限编码依次为 1、2、4、12。用自己的话说明为什么缺失标签不能填均值。
\subsection{C2：只拟合训练统计量}
补 \texttt{fit\_statistics}、\texttt{apply\_statistics} 和 \texttt{to\_tensors}。前者返回训练均值和填补后的总体标准差；第二个只应用已有规则；第三个转换为浮点张量，首列加 1，标签变为列向量。

运行 \texttt{python checkpoints.py C2}。小样例均值为 2，尺度约 0.816497。完成 C1/C2 后运行：
\begin{lstlisting}[language=Python]
import student
train, val, state = support.prepare_train_val(
    student, support.DATA_DIR)
print(train[0].shape, train[1].shape)
print(val[0].shape, val[1].shape)
\end{lstlisting}
预期训练为 \texttt{(8000,9)} 与 \texttt{(8000,1)}，验证为 \texttt{(1000,9)} 与 \texttt{(1000,1)}。训练与验证均无缺失标签；训练就业年限缺失 503 条、债务收入比缺失 1 条，验证就业年限缺失 48 条；测试文件尚未读取。C1 小例子专门练习删除缺失标签。
\begin{questionbox}[验收证据]
列出八列的顺序；解释新增常数列为何放在首列；用数值证明验证值 100 没有参与计算填充值。不要把标签、编号或原始行索引加入特征矩阵。
\end{questionbox}
\clearpage
\section{C3--C4：批次与模型}
\subsection{C3：构造 Dataset 和 DataLoader}
补 \texttt{LoanDataset.\_\_len\_\_}、\texttt{\_\_getitem\_\_} 和 \texttt{make\_loader}。构造函数已经保存特征与标签，不必另做数据复制。
\begin{lstlisting}[language=Python]
dataset = student.LoanDataset(*train)
loader = student.make_loader(dataset, 64, True)
print(len(dataset))
for X_batch, y_batch in loader:
    print(X_batch.shape, y_batch.shape)
\end{lstlisting}
星号 \texttt{*train} 把二元组展开为两个参数，等价于传入 \texttt{train[0], train[1]}。代码应输出 8000 个样本，125 批且每批 64 条。小自检命令为 \texttt{python checkpoints.py C3}，其五条数据应分为 2、2、1。
\subsection{C4：注册参数，计算得分}
补 \texttt{sigmoid}、\texttt{LogisticRegression.\_\_init\_\_} 和 \texttt{forward}。
\texttt{nn.Module} 是可训练模型的基类，\texttt{nn.Parameter} 用来注册权重，使优化器能找到它。输入已经带常数列，因此这里不再增加单独偏置。
\begin{lstlisting}[language=Python]
model = student.LogisticRegression(train[0].shape[1])
logits = model(train[0][:3])
probabilities = student.sigmoid(logits)
print(logits.shape, probabilities.shape)
print(probabilities)
\end{lstlisting}
零初始化时前三条概率都是 0.5；单层逻辑回归不存在隐藏单元的对称性问题，可以这样初始化。
\begin{examplebox}[先手算，再检查]
三个参数为 $(1,2,-1)^\top$，输入为 $(1,3,2)^\top$，得分为 5、概率约为 0.993307。这里 $\top$ 表示转置，括号内是列向量；本例仅有两个非恒定特征，与完整数据的八维输入分开使用。
\end{examplebox}
运行 \texttt{python checkpoints.py C4}。若返回 0.993307 导致得分测试失败，检查是不是在 \texttt{forward} 中过早调用了 sigmoid。若形状是 \texttt{(B,9)}，检查是否把矩阵乘法 \texttt{@} 写成逐元素乘法 \texttt{*}。
\begin{questionbox}[验收证据]
记录首批与末批形状，解释 \texttt{shuffle} 为什么只用于训练；画出“输入矩阵 → 得分列 → 概率列”的形状链。
\end{questionbox}
\clearpage
\section{C5--C6：损失、正则化与更新}
\subsection{C5：确认损失输入与惩罚范围}
补 \texttt{data\_loss} 和模型的 \texttt{penalty}。BCE（Binary Cross-Entropy，二元交叉熵）在本程序通过 \texttt{BCEWithLogitsLoss} 计算，输入为得分，不是概率。\texttt{penalty} 只取 \texttt{theta[1:]}，不含首项截距。

L1（L1 Norm，一范数）惩罚是权重绝对值之和；L2（L2 Norm，二范数）惩罚在本实验明确约定为权重平方和的一半。\texttt{kind="none"} 返回同设备的标量零，训练总目标等于数据损失。

运行 \texttt{python checkpoints.py C5}。得分全为 0 时，平均损失约 0.693147；参数 $(10,3,-4)^\top$ 的 L1 与 L2 惩罚为 7、12.5。自检还检查极端得分的有限损失与惩罚梯度，防止悄悄重复 sigmoid 或误惩罚截距。
\subsection{C6：补全一个可解释的训练循环}
补 \texttt{train\_epoch}。训练器使用 SGD（Stochastic Gradient Descent，随机梯度下降）；本实验每次基于一批样本更新。按以下顺序连通计算：
\begin{enumerate}
\item \texttt{optimizer.zero\_grad()}：清除旧梯度。
\item 用 \texttt{model(X)} 得到得分，计算平均数据损失。
\item 将 \texttt{strength * model.penalty(kind)} 加到数据损失上。
\item 对总目标调用 \texttt{backward()}，再调用 \texttt{optimizer.step()}。
\item 仅在记录时提取普通数值；按每批实际样本数累加数据损失。
\end{enumerate}
运行 \texttt{python checkpoints.py C6}。两条手算样本、零初始化、学习率 0.4、无正则化时，更新后参数应为 $(0,0.1,-0.1)^\top$。它还核验第二步更新，检查是否忘了清梯度。
\begin{pitfallbox}[记录和优化不要混为一谈]
\texttt{train\_epoch} 返回各批更新前数据损失的加权平均，用于理解迭代过程；正式损失曲线由 \texttt{workflow.py} 在每轮结束后，固定同一组参数重新评价整个训练与验证集，二者定义不同。
\end{pitfallbox}
\begin{questionbox}[验收证据]
手写两条样本的梯度；保存第一次更新参数；用五条样本、批量 2 的小例子，解释末批只有 1 条时为何不能把各批平均损失直接等权平均。
\end{questionbox}
\clearpage
\section{C7：选方案，再做最终评价}
\subsection{完成指标，运行验证实验}
补 \texttt{binary\_metrics}，预测概率大于或等于阈值时判为违约。先完成 C7 的小样例检查，再运行全部自检与候选方案训练：
\begin{lstlisting}[language=bash]
python checkpoints.py all
python run_lab.py
\end{lstlisting}
默认比较无正则、L2 系数 0.01、L1 系数 0.01 三组。每组学习率 0.05、最多 60 轮、批量大小 64、种子 42。每次候选训练前重设种子，零初始化权重；比较条件保持一致。

每组保存验证数据损失最小的那一轮，再在三组中选择验证数据损失最小者；保存的是参数副本而非仍随训练变化的引用。固定阈值 0.5。\file{outputs-tianchi-balanced/} 中会出现：
\begin{itemize}
\item \file{candidates.csv}：三组的最优轮数、验证损失和准确率。
\item \file{none_loss.png}、\file{l1_loss.png}、\file{l2_loss.png}：无正则数据损失曲线；同名 history 表保存数值。
\item \file{selection.json}、\file{preprocessing.json}、\file{selected.pt}：冻结的模型设置、预处理统计量与权重。
\item \file{validation_thresholds.csv}：验证集上三种阈值的指标和教学代价；\texttt{example\_cost} 按“漏判代价 2、误报代价 1”计算，仅用于讨论。
\end{itemize}
\subsection{冻结后，单独运行最终测试}
检查上述文件，写下选用哪组、哪轮及理由，然后执行：
\begin{lstlisting}[language=bash]
python run_lab.py --final-test
\end{lstlisting}
这个命令读取已经保存的权重和预处理规则，才首次读取测试文件，输出 \file{test_metrics.json} 与 \file{test_predictions.csv}。请沿用同一数据目录和输出目录。完成后不再基于测试结果重调参数。

TP（True Positive，真正例）、FP（False Positive，假正例）、FN（False Negative，假负例）、TN（True Negative，真负例）均以“违约”为正类。输出包含准确率、精确率、召回率和 F1 分数（F1 Score），并给出常数基线准确率；训练两类相同时，预先约定总预测未违约。分母为 0 的指标约定记为 0。
\begin{questionbox}[验收证据]
展示一张曲线、一张混淆矩阵；说明多数类基线和模型的差别。若准确率尚可但漏判较多，指出应查看的指标，并在验证阈值表中解释误报与漏判如何变化。
\end{questionbox}
\clearpage
\section{实验报告与提交}
\subsection{四部分报告结构}
提交 \file{report.pdf}。PDF（Portable Document Format，可移植文档格式）报告首页写实验名称、学号和姓名；可填写 \file{report-template.md} 后导出，无需同时提交模板原文件。沿用实验 01 的提交方式，在报告中增加分阶段实现与验证。
\begin{enumerate}
\item \textbf{实验概况与完成总览。}简述任务，列出 C0--C7 的完成状态；未完成项写清函数和具体问题。
\item \textbf{C0--C7 实现与验证。}按下页要求逐项写关键实现、小例子的输入及预期与实际输出、结果解释。C0 记录提供的示例及修改观察；C6 保留一次完整参数更新的简短手算与代码对照。
\item \textbf{综合实验结果。}说明数据来源、清洗前后三份样本数和训练统计量处理；列出种子、学习率、批量大小、最大轮数。比较三方案的系数、最优轮数、验证数据损失和准确率；测试前写下选择及理由。插入选定方案的一张训练/验证损失曲线。列出测试准确率、精确率、召回率、F1 分数（F1 Score）、多数类基线和混淆矩阵；矩阵行是真实类别、列是预测类别，0 为未违约、1 为违约。
\item \textbf{结果分析与总结。}解释损失变化、误报与漏判、多数类基线；引用验证阈值表中至少两个阈值，讨论任务代价。不根据测试结果重新调参；说明一项局限及实验收获或建议。
\end{enumerate}
\begin{teacherbox}[把代码与证据对应起来]
关键代码通常摘录约 3--10 行，注明文件与函数；以表达完整为准，不贴整个文件。预期值来自手算或手册，实际值必须自己运行得到；只贴 \texttt{PASS} 不能说明实现。综合图表集中放在报告第 3 部分，各检查点可直接引用。未完成项如实说明，不按篇幅或截图数量评价。
\end{teacherbox}
报告中的图表与数值必须来自本次提交代码，不得把本次抽样称为教材原有划分或照抄教师结果。无需安装证明、完整日志或整篇理论推导。
\subsection{压缩包保留哪些文件}
提交 \texttt{学号\_\allowbreak{}姓名\_\allowbreak{}lab03.zip}，解压后根目录包含 \file{report.pdf} 与 \file{code/}。后者保留 \file{student.py}、\file{run_lab.py}、\file{support.py}、\file{workflow.py}、\file{checkpoints.py}、\file{requirements.txt}，以及 \file{data/classroom/} 中三份数据、划分清单与数据说明；不打包天池原始大文件。

采用 Notebook 入口时另附 \file{logic.ipynb}，仍须保留学生脚本和辅助模块。与实验 01 不同，当前训练入口会调用检查模块，所以 \file{checkpoints.py} 也须随包提交。完整目录示意见本实验 \file{README.md}。

输出目录中的原始表、权重与预测明细自行保留即可；无需重复提交报告中已有的图表。不打包 Python 环境、安装包、缓存、教师答案与课程讲义。截止时间和提交平台以课程通知为准。
\clearpage
\subsection{每个检查点在报告中写什么}
\par\noindent\begin{tabularx}{\linewidth}{L{1cm}L{4.4cm}Y}\toprule
阶段 & 实现定位或观察对象 & 至少记录的一项数值验证及解释\\\midrule
C0 & 提供的矩阵乘法示例；运行入口与版本 & 输入、权重与输出形状；改动一个权重前后的预期及实际得分。说明变化原因。\\
C1 & \file{filter_labels}、\file{encode_employment} & 含缺失标签的小表保留多少行；特殊年限、普通年限与缺失值的编码。解释标签与输入的区别。\\
C2 & \file{fit_statistics}、\file{apply_statistics}、\file{to_tensors} & 小例子的均值、填补后标准差、验证转换值与张量形状/类型；改变验证数据后训练统计量不变。\\
C3 & \file{LoanDataset} 的两个方法、\file{make_loader} & 同索引的特征和标签；样本数不能整除批量大小时的各批大小与总数。说明打乱规则。\\
C4 & \file{sigmoid}、模型初始化与 \file{forward} & 一组输入和参数对应的得分、概率与形状。解释参数注册，以及得分和概率的区别。\\
C5 & \file{data_loss}、\file{penalty} & 零得分损失；一组参数的三种惩罚。只改截距后惩罚不变；区分数据损失与总目标。\\
C6 & \file{train_epoch} 的五步更新 & 两样本例子的初始参数、概率、梯度、学习率及一次更新的手算/运行对照；第二次更新结果。解释清梯度与按样本数加权。\\
C7 & \file{binary_metrics}、提供的选择与测试流程 & 小例子在两个阈值下的类别、混淆矩阵和指标，含概率等于阈值的情况。说明验证选择与冻结测试的职责。\\\bottomrule
\end{tabularx}
\begin{examplebox}[一条记录应能被复核]
例如 C2：写出自己选取的训练列与验证值，摘录拟合和复用统计量的代码；列出预期与实际均值、尺度和转换值，最后解释验证值为何不能参与拟合。不要只写“标准化完成”或“C2 通过”。具体函数和填写位置见报告模板第 2 节。
\end{examplebox}
\clearpage
\section{提交前检查、评分与拓展}
\subsection{检查可复现性，再打包}
\begin{enumerate}
\item 运行 \texttt{python checkpoints.py all} 核对实现，不要求验证截图。
\item 确认报告来自待提交代码的训练、验证选择与冻结后最终测试；测试结果揭晓后不再据此调参。
\item 打开 \file{report.pdf} 检查文字、曲线与数值。解压到另一目录，运行 C0 和全部自检，确认代码与数据齐全；不必为打包重复测试选模。
\end{enumerate}
公开检查帮助定位问题，不等于全部评分规则。尚未解决的问题在报告中如实注明，保留已完成实现。
\subsection{评分结构对齐}
\par\noindent\begin{tabularx}{\linewidth}{L{2cm}Y L{2.2cm}}\toprule
比例 & 验收内容 & 对应阶段\\\midrule
30\% & 环境 5\%、读表 5\%、缺失值与训练统计量 10\%、数据加载 10\% & C0--C3\\
40\% & 参数初始化 10\%、sigmoid 10\%、前向传播 10\%、正则化 10\% & C4--C5\\
30\% & 损失 10\%、训练及冻结后的最终评价 10\%、收获与建议 10\% & C5--C7\\\bottomrule
\end{tabularx}
此建议沿用指定指导书第 1.6 节的 30/40/30 结构，按 checkpoint 核验实现与防泄漏要求。没有“必须达到某个准确率”的硬门槛。\texttt{PASS} 不替代理解；能说明尚未解决的问题也是有效实验记录。
\subsection{课后问题与参考要点}
\begin{enumerate}
\item 缺失标签能否按多数类补齐？不能，伪造的监督信号会改变训练与评估含义。
\item 预测全为未违约为什么可能准确率很高？类别不平衡；同时看正类召回率与多数类基线。
\item 过大正则化一定改善泛化吗？不会；它可能抑制必要的信号，须用验证集比较。
\item 为什么本实验不要求手写反向传播？先建立损失、计算图与更新的对应；后续实验再实现和比较梯度算法。
\end{enumerate}
\subsection{课堂数据与自备划分}
默认从 \file{code/data/classroom/} 读取真实贷款子集，结果写入 \file{code/outputs-tianchi-balanced/}，不能复用旧模拟权重。天池 \file{testA.csv} 没有标签，提交示例中的 0.5 也不是真值；课堂三份数据均从带标签训练文件中重新划分。

自备划分目录须含 \file{train.csv}、\file{val.csv}、\file{test.csv}，满足八字段与真实标签约定。\file{dti} 保留原尺度，不直接替换成模拟负债比例。更多字段不会自动进入模型，须核对类型、单位与预测时可用性后再扩展。
\begin{lstlisting}[language=bash]
python run_lab.py --data-dir ../course-data --output outputs-course
python run_lab.py --data-dir ../course-data --output outputs-course --final-test
\end{lstlisting}
\clearpage
\section{排错、结果解读与资料对照}
\subsection{常见问题按现象定位}
\par\noindent\begin{tabularx}{\linewidth}{L{3.5cm}Y}\toprule
现象 & 首先检查\\\midrule
找不到 \texttt{torch} & 终端解释器是否为课程环境；Notebook 内核是否一致\\
\texttt{NoneType} 报错 & 当前 checkpoint 的占位返回是否已替换，是否漏补依赖阶段\\
输入/标签形状不匹配 & 标签是否为 \texttt{(B,1)}；是否少了增广常数列\\
所有概率都相同 & 参数是否注册为 \texttt{nn.Parameter}；训练是否执行更新\\
损失不降或出现非数值 & 特征是否填充且标准化；学习率是否过大；损失是否接收原始得分\\
第二步手算不一致 & 是否每批清梯度；是否对总目标而非记录值反向传播\\
预测与标签错位 & 是否独立打乱了标签；Dataset 是否按相同索引返回二者\\
验证很高、测试很差 & 是否混用数据、把标签放入特征、根据测试结果选参数\\\bottomrule
\end{tabularx}
\subsection{参考运行仅用于核对流程}
本次真实贷款子集使用 8000/1000/1000 条记录；按指导书将训练正负各设为 4000 条，验证与测试各含 200 条违约。零初始化、训练种子 42、学习率 0.05、批量 64，每组最多 60 轮。按验证数据损失选择 L1 系数 0.01 的第 3 轮，阈值固定为 0.5。

验证损失约 0.648418、准确率 0.644、召回率 0.65；测试损失约 0.620095、准确率 0.667、精确率约 0.320、召回率 0.59，混淆矩阵为真负例 549、假正例 251、假负例 82、真正例 118。常数基线准确率为 0.8。

原自然比例训练的测试召回率为 6\%，平衡后为 59\%，但误报由 3 增至 251。\textbf{这是同一已查看测试集上的教学对照，不是新的盲测，也不是教材原结果。}验证阈值 0.3/0.5/0.7 时召回率为 95\%/65\%/19\%，示例代价为 699/426/373；不能仅凭召回率判断优劣。平衡训练改变类别先验，输出未经概率校准。不同库版本可能影响结果，不按逐位相等评分。
\subsection{资料与改编对应}
指定《人工智能数学原理与算法：实验指导书（最终）.pdf》位于项目上级目录。第 1.4 节对应 C1--C2；第 1.5.2 节对应 C3--C6；第 1.5.3 节与第 1.6 节对应 C7、提交及评分建议。页码为书内第 1--13 页，文件第 15--27 页。

仓库 \file{source/experiment/Logic.ipynb}、\file{source/experiment/utils.py} 用于核对旧版入口与字段转换；前置讲义列出理论课件与课程总表路径。随包 \file{code/data/README.md} 与 \file{teacher/prepare_classroom_data.py} 给出真实数据来源及抽样规则；教师实现与验证说明位于 \file{teacher/}，学生下载包不包含该目录。报告提交方式参考同级实验 01 的实验手册第 7 节及报告模板；实验 02 不另交报告。分阶段记录方式借鉴用户提供的《数据科学基础》个人 HW1 成稿第 5、9--10 页的“任务、代码、输出、结论”组织；只借鉴结构，外部资料位置见教师说明。

相较原版，明确得分/概率接口，保存训练统计量，处理 \texttt{1 year} 单数形式，并用参数首项表示截距；零初始化和三组预设比较是课堂改编。数学符号完整定义及分类索引见《前置讲义》第 10 节。

\end{document}
